\label{chap:conclusions}
This work, named "ColonVision: Automated Detection of Colorectal Anomalies via Computer Vision", presents a survey of multiple \gls{DL} architectures for \gls{CRC} classification. Multiple different approaches were taken, including training models with different \gls{DS} and validating their output with \gls{CL} images from other sources and comparing the results with many optimization techniques. In addition, it was tested the performance of a hybrid quantum-classical model, which transform raw medical images into a feature-rich representation, enhancing classification performance. Despite the many differents scenarios tested, some of the used \gls{DS}, such as \gls{LC25000} and \gls{CRC5000}, were unbalanced and had different normalization techniques applied, what may have influenced in the results in chapter \ref{chap:results}.

The results exposed in chapter \ref{chap:results} answer some of the questions introduced in chapter \ref{chap:related} with other researchers works. Classic \gls{DL} models can be used for the detection and classification of colorectal anomalies with high reliability in predictions. Models with less parameters such as AlexNet had promising performance as well when trained with \gls{NCTCRCHE100K} with no previous normalization image set, with a higher accuracy and f-score higher than ninety percent using \gls{SGD} and Adammax optimizers.

\gls{QNN} models demonstrated a competitive robustness and promising results even when compared to more complex \gls{CNN} and \gls{RNN}. However, the large train time, specially for higher resolution images, still an area of improvement to enable more in-depth studies. Also, for Google colab, larger \gls{DS} are narrowed by the time limit of continuous execution. The improvement in python libraries of quantum hardware simulation and cloud \gls{ML} service support for quantum computation may spread the architecture viability in multiple different areas, including tumor classification.

When considering the different \gls{DS} proposed in chapter \ref{chap:methodology}, \gls{NCTCRCHE100K} image set with no previous normalization, NCT-CRC-HE-100K-NONORM, was the best one to train the \gls{DL} models for generalization with images from other sources. Although the detection of images from \gls{LC25000} lack of higher performance, the absence of the Macenko colorization technique \cite{5193250} limits the conclusiveness of this results.

There are multiple improvements points in the study, which could be explored in future works. The implementation of Macenko colorization technique \cite{5193250} in \gls{LC25000} to validate the trained architectures with \gls{NCTCRCHE100K} performance. Furthermore, the train with \gls{CRC5000} with augmentation techniques is a great option to analyse its attainment in being used as base image set for training. Finally, modifications in the proposed hybrid quantum-classical models \cite{quanvolution} presents lots of room for improvement to enhance robustness of the models. The real relevance of this survey would be achieve by combining the proposed architectures with resource efficiency techniques to deploy those models on edge devices, democratizing this classification technology to the major population.